\documentclass[11pt]{article}
% ---- style ----
\usepackage[utf8]{inputenc}
\usepackage[T1]{fontenc}
\usepackage{lmodern}
\usepackage{microtype}
\usepackage{amssymb}
\usepackage[a4paper,margin=0.85in]{geometry}
\usepackage{xcolor}
\definecolor{Accent}{HTML}{1565C0}
\usepackage{titlesec}
\titleformat{\section}{\normalsize\bfseries\color{Accent}}{}{0em}{}
\titlespacing*{\section}{0pt}{6pt}{2pt}
\usepackage{enumitem}
\setlist{nosep,leftmargin=1.2em}
\usepackage[colorlinks=true,urlcolor=Accent,linkcolor=Accent]{hyperref}
\input{glyphtounicode}\pdfgentounicode=1
\setlength{\parindent}{0pt}
\setlength{\parskip}{4pt}
\pagestyle{plain}
% ---------------
\begin{document}

\begin{flushright}
Alexander Salganik\\
Department of Physics and Astronomy, University of Turku, Finland\\
\href{mailto:alexander.a.salganik@utu.fi}{alexander.a.salganik@utu.fi} $\cdot$
\href{https://salganik.me}{salganik.me}\\
\today
\end{flushright}

Fellowship Selection Committee\\
MIT Kavli Institute for Astrophysics and Space Research\\
Massachusetts Institute of Technology

Dear Members of the Selection Committee,

I am applying for the MIT Kavli Postdoctoral Fellowship and the Juan Carlos Torres
Fellowship in Astrophysics. I am a third-year PhD candidate at the University of
Turku, working with Prof.~Juri Poutanen and Dr.~Sergey Tsygankov, and I will defend
my thesis in summer 2027, before the start of the fellowship.

My research asks how matter accretes onto strongly magnetized neutron stars, and I
answer it by finding the right observation at the right moment. Most of my results come
from campaigns I proposed myself, often as target-of-opportunity requests during
unexpected outbursts. This is how I obtained the first X-ray polarization detection from
a long-period pulsar and how I caught the frequency hysteresis of quasi-periodic
oscillations in 1A~1118$-$61, a sign that the accretion flow remembers its history. My
CV and research statement give the details; here I want to explain why MKI is the place
where this work can grow.

My proposal links three diagnostics that are usually studied separately: cyclotron lines
(field strength), X-ray polarization (field orientation and emission geometry) and timing
(how the flow evolves). The natural next step is a band no instrument has yet opened for
accreting pulsars. GOSoX, selected by NASA for launch in 2030, will measure polarization
at 0.2--0.4\,keV, where many pulsars show a soft excess produced by reprocessing on the
inner disc. Predicting and interpreting this signal --- starting with Her~X-1, whose
geometry IXPE has already measured --- is a project that only makes sense at MKI, and the
fellowship years coincide with the mission's preparation. I would be excited to work with
Herman Marshall's group on the GOSoX science case, and to connect the timing part of my
programme with the X-ray timing expertise of Erin Kara and Deepto Chakrabarty.
% TODO: adjust names after contacting people; if someone replied, say so explicitly.

In return, I would bring experience that complements MKI's strengths: hands-on IXPE
spectro-polarimetric analysis, a track record of rapidly obtaining NuSTAR, IXPE, NICER
and Swift time for transients, ongoing collaborations with the IXPE pulsar community,
including members of the GOSoX team in T\"ubingen, and access to dense SRG/ART-XC and
Insight-HXMT monitoring through my collaborators.

Mentoring matters to me. I have built my career across two countries and in
collaborations spanning Finland, Russia, China, Germany and the UK, and I know how much
early guidance and a welcoming group matter, especially for international students. I
currently supervise student research projects in the Cambridge Future Scholar
Programme, guiding students from a first research question to a written paper, and I
taught practical X-ray data analysis to MSc students in Turku. As a Kavli Fellow I would
mentor undergraduate researchers, for example through MIT's UROP programme, and help
make the group a place where every student feels they belong.

My application includes my CV, publication list (NASA ADS), research statement and
research proposal. Letters of recommendation will be provided by Prof.~Juri Poutanen,
Dr.~Sergey Tsygankov and [third referee]. Thank you for your consideration.

Sincerely,\\[3pt]
Alexander Salganik

\end{document}
